\chapter{Interventions at the Cut}
\label{ch:do}

\section{What is being built, and why here}

A program in rho with graded where-clauses over $\Rnn$ already carries a causal model. The
likelihood of a causal history is a product over firing sites of a locally normalized factor
(Theorem~\ref{sci:thm:score} will take its derivative); each site's factor is its clause; each
site's parents are the data its pattern and clause read; and the resolver that settles its race
is the exogenous noise. This chapter makes the model explicit and adds the one thing Pearl's
calculus of causation has that the clause language lacked: an operator that \emph{intervenes}.
It follows the working note in which the construction was first set out~\cite{Interventions}.

The operator is a scoped decoration on parallel composition,
\[
  \dok(\,y_1 \models \varphi_1 \ \&\ \cdots\ \&\ y_n \models \varphi_n\,)\ \ink\ P ,
\]
read as: whenever a race for the binding of $y_i$ arises inside $P$, resolve it among the
candidates satisfying $\varphi_i$, with the site's own weights restricted to that set. It is placed
at the cut because an aperture is a property of the cut and of neither agent flanking it
(Remark~\ref{gw:rmk:noncompositional}), and because whether a particular cut will ever occur is
undecidable even in the sequential fragment, so an intervention cannot be a surgery on a known node
and must be a policy conditional on an event. The construct lowers into the where-clause by the
normal-form theorem (Theorem~\ref{gw:thm:dnf}), so it is a conservative extension, and it sits
beside the cost membrane of Chapter~\ref{ch:carho} as a second index on a polymorphic cut.

With it, Pearl's causal Bayesian networks over finite domains embed exactly: a variable is a join on
its parents' channels together with an ancilla race over its domain, the conditional probability
table is the clause in disjunctive normal form, and the truncated factorization is the likelihood
of a causal history. Intervening on the ancilla binding deletes the incoming edges on the nose, the
three rules of the do-calculus are theorems about the embedded fragment rather than axioms, and
counterfactuals are replay from a recorded pre-state under the intervention.

The chapter sits here, directly after Chapter~\ref{ch:grading}, because it is the first thing one
builds with a graded clause that one could not have built with a weight map, and because two parts
of the book downstream need it: the mortal scientist of \chSci, whose hypotheses become
rung-two objects once they can be placed under a $\dok$ (\S\ref{do:sec:scientist}), and the
Prestige's argument that a proof system or a learner that cannot inhabit a race cannot do what
evolved intelligence does --- for which this construct is the constructive content. Pearl himself
has not appeared in this book before now, so the ladder is introduced before it is climbed.

\section{Three rungs}
\label{do:sec:ladder}

Pearl's ladder of causation~\cite{Pearl09} has three rungs, and a learner's power is where it
stands on it. \emph{Association}: what one can learn from a record of what happened ---
$\Pr(y \mid x)$, the distribution of completed runs. \emph{Intervention}: what one can learn by
acting --- $\Pr(y \mid \dok(x))$, the distribution when $X$ is \emph{set} to $x$ rather than
observed to be $x$. \emph{Counterfactual}: what would have happened in the run that actually
occurred, had $X$ been otherwise --- $Y_x(u)$, with $u$ the particular draw of everything the
model does not explain. Each rung is strictly stronger than the one below, in the sense that
there are pairs of models agreeing at one rung and disagreeing at the next.

\begin{definition}[Causal Bayesian network]
\label{do:def:cbn}
A \emph{causal Bayesian network} $M = (G,\{P_i\})$ over variables $V_1,\dots,V_n$ with finite
domains $D_1,\dots,D_n$ consists of a DAG $G$ on the $V_i$ and, for each $i$, a conditional
probability table $P_i(v \mid \mathit{pa})$ for $v \in D_i$ and
$\mathit{pa} \in \prod_{j\in\mathrm{PA}_i} D_j$. It induces
$\Pr_M(v_1,\dots,v_n) = \prod_i P_i(v_i \mid \mathit{pa}_i)$. $M$ is \emph{positive} when every
table entry is strictly positive. For $X \subseteq V$ and an assignment $x$, the \emph{submodel}
$M_x$ replaces $P_i$ by the point mass at $x_i$ for each $V_i \in X$ and deletes the incoming
edges of $X$, and $\Pr_M(y \mid \dok(x)) := \Pr_{M_x}(y)$.
\end{definition}

A structural causal model refines this by writing each $V_i = f_i(\mathrm{PA}_i, U_i)$ with
$U_i$ independent exogenous variables; a CBN is a structural model with the $U_i$ marginalized,
and \S\ref{do:sec:cf} recovers the $U_i$ as the record of the resolver's draws.

The reader who has followed the book this far will notice that the ladder is a statement about
\emph{time}. Association is what a trace shows; a trace is linear time. Intervention is what an
experimenter who supplies contexts can learn; that is branching time, and context-labeled
bisimulation is its equivalence. Counterfactual is replay with the record held. The alignment is
made exact below, and it is the reason this chapter is in the Machinery part and not in a chapter
on statistics.

\section{Interventions have one carrier}
\label{do:sec:carrier}

Before adding syntax it is worth settling what an intervention can touch, since the answer
constrains the syntax.

\begin{definition}[Intervention, semantically]
\label{do:def:intervention}
An \emph{intervention} at a receipt site $r$ of $P$ is a replacement of the map
$\Cand(P,r) \to \Semi$ denoted by $\psi_r$ with another such map. It is \emph{hard} when the
replacement is the indicator of a single candidate, \emph{set-valued} when it is the indicator of
a subset, and \emph{soft} otherwise.
\end{definition}

\begin{proposition}[Deterministic communication is intervention-transparent]
\label{do:prop:transparent}
At a site with exactly one candidate, every intervention is either the identity on the normalized
dynamics or disables the site. In particular no intervention at such a site changes which outcome
is reached when the site fires.
\end{proposition}

\begin{proof}
The fibre is a singleton; a map on a singleton is determined by whether its value is zero.
Normalization within the contention component sends any nonzero value to $1$.
\end{proof}

To intervene at a deterministic site one must change the fibre --- inject a datum or delete one ---
and the moment a datum is injected the site has two candidates and is a race. The carrier of
intervention is therefore the candidate set at a receipt, and a nontrivial intervention exists
exactly where that set has, or is given, more than one element.

\begin{proposition}[No hard intervention by context over $\Rnn$]
\label{do:prop:nocontext}
Let $r$ be a race in $P$. No rho context $C[-]$ makes $C[P]$ resolve $r$ to a chosen candidate
with probability $1$ under the normalized $\Rnn$ semantics, unless $C$ consumes every competing
candidate before $r$ can fire --- which is itself a race $C$ wins with probability strictly less
than $1$.
\end{proposition}

\begin{proof}
This is Principle~\ref{gw:prin:aperture} read in the forcing direction. A context can add
contention and consume contenders, and consuming a contender is a candidate at a site in
contention with $r$. Over $\Rnn$ every candidate of positive weight has positive probability
(Corollary~\ref{gw:cor:adequacy} with positive constants), so $r$ fires first with positive
probability.
\end{proof}

So over $\Rnn$ a hard intervention requires access to the \emph{clause}: one must be able to
write at the site, not merely around it. This is the observation/intervention gap in its
sharpest form --- what a context can do is observe and perturb; what an owner of the code can do
is set --- and it is why the construct below is a decoration the programmer writes rather than a
context one composes. One qualification: the clause controls which candidate fires and not when.
An intervention on timing is outside the clause, and outside this chapter, unless time is reified
as data --- at which point it is a race again.

\begin{principle}[Interventions are metered because they are races]
\label{do:prin:metered}
A race is a COMM and a COMM is charged. Reading the trace is free; intervening costs. This is not a
policy choice: intervention has no carrier other than the event the ledger already prices.
\end{principle}

\section{The construct: a decorated cut}
\label{do:sec:construct}

There is always a node in the syntax tree above a race, and it is a parallel composition: a race
needs a receiver and at least two things it could meet. The cost membrane of
Chapter~\ref{ch:carho} already made the cut polymorphic, indexing it by seals; an intervention is a
second index on the same operator. And it cannot be anything more specific than a scoped
decoration, for a halting reason: whether a particular cut will ever occur is undecidable, already
in the sequential fragment, since ``this receipt fires'' is a nontrivial reachability property. So
an intervention cannot be a surgery on a node known in advance. It is a \emph{policy conditional on
an event} --- if a race for this binding arises, resolve it thus --- vacuous if the race never
arises. Pearl's nearest object is a dynamic treatment regime rather than an atomic $\dok$.

\begin{definition}[The $\dok$ membrane]
\label{do:def:do}
Extend the syntax of Chapter~\ref{ch:grading} with
\[
  P ::= \cdots \;\mid\; \dok(\,C\,)\ \ink\ P, \qquad
  C ::= y \models \varphi \;\mid\; y \models \psi \;\mid\; C \,\&\, C ,
\]
where $y$ is a pattern variable, $\varphi$ a crisp formula and $\psi$ a graded formula, each with
$y$ among its free names. Write $\dok(y = Q)$ for $\dok(y \models \crisp{y = @Q})$ and
$\dok(y \in S)$ for the disjunction over a finite set $S$.
\end{definition}

\begin{lstlisting}
do( y = Q2 ) in { x!( Q1 ) | for( y <- x )P | x!( Q2 ) }

do( y1 |= phi1 & y2 |= phi2 ) in { P1 | P2 | P3 }
\end{lstlisting}

\begin{definition}[Membrane semantics]
\label{do:def:membrane}
Let $r$ be a receipt occurrence inside the scope of $\dok(C)\ \ink\ P$ binding a variable $y$
constrained in $C$ by $\varphi$ (resp.\ $\psi$). At $r$ the clause is read as
$\psi_r \tensorv \crisp{\varphi}$ (resp.\ $\psi_r \tensorv \psi$), with the constraint evaluated
under the candidate's substitution, and normalization within the contention component is taken
\emph{after} the restriction. Reduction is $\dok(C)\ \ink\ P \longrightarrow_w \dok(C)\ \ink\ P'$
whenever $P \longrightarrow_w P'$ under clauses so read. A receipt in scope binding an unconstrained
variable is unaffected; a receipt arriving in scope by reflection is constrained like any other.
\end{definition}

\begin{definition}[Lowering]
\label{do:def:lowering}
On a scope $P$ in which every receipt binding a constrained variable is syntactically present and
its channel is a literal or \textsf{new}-bound name, let $\lfloor \dok(C)\ \ink\ P \rfloor$ be $P$
with each such receipt's clause replaced by $\psi_r \tensorv \crisp{\varphi}$ (resp.\
$\psi_r \tensorv \psi$) and the $\dok$ erased.
\end{definition}

\begin{proposition}[Lowering is exact on flat scopes]
\label{do:prop:lowering}
If $P$ contains no drop $*z$ at a position that can come to contain a receipt, then
$\dok(C)\ \ink\ P$ and $\lfloor \dok(C)\ \ink\ P \rfloor$ have the same weighted transition system,
hence the same graded denotation and the same likelihood on every causal history.
\end{proposition}

\begin{proof}
Under the hypothesis every receipt that will ever be in scope is present in $P$ up to substitution
of names for pattern variables, which does not create or destroy binders. The membrane reads each
such clause exactly as the lowering rewrites it, and the contention relation is unchanged because
restriction removes candidates without adding any.
\end{proof}

The gap between the two semantics is exactly the dynamic scope that reflection introduces: received
code is the one route by which an unlisted binder --- and, in the causal reading, an unlisted
parent --- enters. We take the membrane as the definition and the lowering as the optimization.

\begin{observation}[Restrict, not override]
\label{do:obs:restrict}
The membrane \emph{multiplies} the clause by the constraint; it does not replace it. For
$\fork(y \leftarrow x \;\where\; \crisp{*y > 5})P$ under $\dok(y = 3)$, the restricted clause is
$0$ on every candidate and the site waits; an override would force a $3$. Pearl's atomic $\dok$
severs the mechanism and so is an override in spirit --- but \S\ref{do:sec:pearl} shows that on the
embedded models the two coincide, because the embedding puts the mechanism in the \emph{weights}
and the restriction is on an ancilla whose own gate is $\top$.
\end{observation}

\begin{proposition}[A compressed $\dok$ is a weight map; composition]
\label{do:prop:dnf}
A constraint $y \models \bigoplus_i v_i \tensorv \crisp{\chi_i}$ with the $\chi_i$ pairwise
exclusive denotes a weight map on the fibre at $y$'s site, and every weight map on that fibre is
so denoted (Theorem~\ref{gw:thm:dnf} applied to the constraint). Moreover
$\dok(C_1)\ \ink\ \dok(C_2)\ \ink\ P$ and $\dok(C_1 \,\&\, C_2)\ \ink\ P$ have the same weighted
transition system, $\&$ is commutative and associative, and $\dok(y \models \top)\ \ink\ P$ has the
transition system of $P$.
\end{proposition}

So several interventions compress into one expression for the reason weight maps have a normal
form: a formula has one value. The $\dok$ is a scoped weight map placed at the cut, and the
normal-form theorem is what makes it a conservative extension. It is also, in the vocabulary
\chSci\ will use, the operator that takes a \emph{specification} --- the $\nu$-property ``every
binding of $y$ satisfies $\varphi$'', which one would check --- and installs it in \emph{guard}
position at every site in scope, including sites not yet visible.

\section{The embedding of Pearl's calculus}
\label{do:sec:pearl}

\subsection{The variable gadget}

For a variable $V_i$ with children $V_{j_1},\dots,V_{j_k}$, allocate fresh channels
$v_i^{(j_1)},\dots,v_i^{(j_k)}$, one per child, and a readout channel $v_i^{\circ}$; emitting
$V_i$'s value $d$ is $v_i^{(j_1)}!(d) \mid \cdots \mid v_i^{(j_k)}!(d) \mid v_i^{\circ}!(d)$.

\begin{definition}[The gadget]
\label{do:def:gadget}
For each $V_i$ with parents $V_{p_1},\dots,V_{p_m}$ and domain $D_i = \{d_1,\dots,d_s\}$,
allocate a fresh ancilla channel $w_i$ and write
\begin{lstlisting}
new wi in {
    wi!(d1) | ... | wi!(ds)
  | for( p1 <- vp1_i & ... & pm <- vpm_i
       & u1 <- wi & ... & us <- wi
         where  [ *p1 = pa1 and ... and *pm = pam and *u1 = v ] * P_i(v | pa)
              + ...   // one summand per (pa, v) in the table
       )
    { emit_i( *u1 ) }
}
\end{lstlisting}
where \texttt{vpj\_i} is the copy of parent $j$ allocated to this gadget, the crisp bracket is the
indicator embedding, and the sum has one summand per $(\mathit{pa}, v)$. The \emph{model term}
$\sem{M}$ is the parallel composition of the gadgets under one \textsf{new} binding all the
channels.
\end{definition}

The join consumes all $s$ ancilla messages, so there is no residual on $w_i$ and no which-path
record; the candidates are the $s!$ bijections from $u_1,\dots,u_s$ to $D_i$, and the continuation
mentions only $u_1$, so the $s!$ candidates fall into $s$ outcome classes of $(s-1)!$ each. The
clause is constant on each class. Under normalization within the contention component the
$(s-1)!$ multiplicity cancels, and the probability that the gadget emits $v$ given parent data
$\mathit{pa}$ is $P_i(v \mid \mathit{pa})$.

\begin{lemma}[Staging]
\label{do:lem:staging}
In $\sem{M}$ every receipt is its own contention component, each gadget fires exactly once, and the
causal order of every history is a linear extension of $G$.
\end{lemma}

\begin{proof}
Every channel has exactly one receiver, fan-out giving each child its own copy, and every ancilla
channel is private to its gadget, so no two receipts share a message or a receipt. A gadget is
enabled exactly when all its parent copies have been emitted.
\end{proof}

\begin{theorem}[Adequacy of the embedding]
\label{do:thm:adequacy}
Causal histories of $\sem{M}$ reaching a normal form are in bijection with assignments
$(v_1,\dots,v_n) \in \prod_i D_i$, read off the readout channels, and the likelihood of the
history for $(v_1,\dots,v_n)$ is $\Pr_M(v_1,\dots,v_n)$.
\end{theorem}

\begin{proof}
By Lemma~\ref{do:lem:staging} a history is determined, up to the interleaving the causal quotient
removes, by the outcome class chosen at each gadget, i.e.\ by the vector of values; conversely
every vector is realized since every ancilla value is present. The factor contributed by gadget $i$
is its restricted-and-normalized clause value, $P_i(v_i \mid \mathit{pa}_i)$, and the likelihood
multiplies the factors.
\end{proof}

The truncated-factorization \emph{form} is not imposed on the semantics; it \emph{is} the
semantics. The staging that makes each gadget its own contention component is the condition under
which local probabilistic choice composes to a global measure at all~\cite{VaraccaVolzerWinskel}.

\begin{remark}[The gadget is a stochastic matrix, and it has a quantum twin]
\label{do:rmk:shor}
Definition~\ref{do:def:gadget} is the one-qubit gadget of the complex instance
(Chapter~\ref{sec:quantum}, \cite{GradedQuantum}) --- ancilla channel carrying the domain, a join
that consumes everything, a continuation that re-emits without mentioning what was consumed, a
clause in disjunctive normal form indexed by (input, output) --- with three changes: the domain is
$D_i$ rather than $\{0,1\}$, the inputs are several parents rather than one wire, and the matrix
the clause realizes is \emph{stochastic} rather than unitary. The resemblance is exact and is a
statement about the semiring: the quantum gadget realizes every $2\times 2$ matrix over $\CC$
because the clause may mention both the input and the ancilla; this one realizes every stochastic
matrix over $\Rnn$ for the same reason, with stochasticity supplied by normalization within the
component rather than by the clause. What differs is what one then does with the matrix: over
$\CC$ one arranges for columns to interfere and reads a detector that never fires; over $\Rnn$ one
reads the trace.
\end{remark}

\subsection{Surgery}

\begin{theorem}[Graph surgery is restriction of the ancilla race]
\label{do:thm:surgery}
Let $M$ be positive, $V_i \in V$, $v \in D_i$. Then
\[
  \sem{\,\dok(u_1 = v)\ \ink\ \sem{M}\,} \;=\; \sem{M_{V_i = v}}
\]
as weighted transition systems up to structural congruence, where $u_1$ is the first ancilla
pattern of gadget $i$. Hence $\Pr_{\sem{M}}(y \mid \dok(V_i = v)) = \Pr_{M}(y \mid \dok(V_i = v))$
for every $y$.
\end{theorem}

\begin{proof}
At gadget $i$ the restricted clause is the original clause times $\crisp{*u_1 = v}$, which is
$P_i(v \mid \mathit{pa})$ on the $(s-1)!$ candidates in the class $u_1 \mapsto v$ and $0$
elsewhere. The contention component is the gadget alone, so normalization divides by
$(s-1)!\,P_i(v \mid \mathit{pa})$, nonzero by positivity, and the class fires with probability $1$
\emph{for every value of the parents}. The gadget's normalized dynamics is therefore the point mass
at $v$ independent of $\mathit{pa}$, which is exactly the mechanism of $M_{V_i=v}$ at $V_i$; the
parents are still consumed, so the term is unchanged elsewhere; and every other gadget is unaffected
by Proposition~\ref{do:prop:lowering}, the model term being flat. Theorem~\ref{do:thm:adequacy}
applied to $M_{V_i=v}$ gives the distribution.
\end{proof}

The theorem is the reason to address the \emph{ancilla} binding. Restricting a parent pattern
instead would select on inputs --- a conditioning --- not sever a mechanism. The embedding fixes
that the value of a variable is carried by its ancilla binding, and that is the whole of what makes
the surgery exact: the gadget's own gate is $\top$, so restriction and override coincide, and the
weights that encoded the mechanism are the thing normalization discards.

\begin{remark}[Positivity is starvation]
\label{do:rmk:positivity}
If $P_i(v \mid \mathit{pa}) = 0$ for some reachable $\mathit{pa}$, the restricted gadget has clause
$0$ on every candidate when those parents arrive, and the receipt waits forever: the history
deadlocks by starvation. This is exactly where Pearl's $\Pr(y \mid \dok(x))$ is undefined without
positivity, and the embedding renders the undefinedness as a mortal computation's death rather than
as a division by zero. The reader of \chSci\ will recognize the outcome $\bot$ of
Proposition~\ref{sci:prop:trichotomy}: the environment did not answer.
\end{remark}

\begin{corollary}[The do-calculus is derived]
\label{do:cor:rules}
For positive $M$ and disjoint $X,Y,Z,W \subseteq V$, the three rules of the
do-calculus~\cite{Pearl95} hold of $\sem{M}$ under $\dok$ as defined here:
\begin{enumerate}[leftmargin=1.6em,itemsep=1pt]
\item $\Pr(y \mid \dok(x), z, w) = \Pr(y \mid \dok(x), w)$ if
      $(Y \dsep Z \mid X, W)_{G_{\overline{X}}}$;
\item $\Pr(y \mid \dok(x), \dok(z), w) = \Pr(y \mid \dok(x), z, w)$ if
      $(Y \dsep Z \mid X, W)_{G_{\overline{X}\underline{Z}}}$;
\item $\Pr(y \mid \dok(x), \dok(z), w) = \Pr(y \mid \dok(x), w)$ if
      $(Y \dsep Z \mid X, W)_{G_{\overline{X}\,\overline{Z(W)}}}$,
\end{enumerate}
with conditioning read as post-selection on histories and $\dok$ on several variables as the
conjunction $\&$ of their ancilla restrictions.
\end{corollary}

\begin{proof}
By Theorems~\ref{do:thm:adequacy} and~\ref{do:thm:surgery} and
Proposition~\ref{do:prop:dnf}, every interventional distribution of $\sem{M}$ is the
corresponding distribution of a submodel of $M$, and post-selection on histories is conditioning
in that submodel. The rules are Pearl's theorem for the CBN $M$~\cite{Pearl95,Pearl09}.
\end{proof}

The content of the corollary is not that the rules are true --- Pearl proved that --- but that the
object they are stated for is literally present in the semantics, so that nothing is being
modeled. The d-separation premises are statements about $G$, and $G$ is the causal order of the
staged term; whether they are expressible as crisp-sort formulae over $\sem{M}$, and so
dischargeable by the same checker that discharges everything else, is the first of the open items
in \S\ref{do:sec:open}.

\subsection{Counterfactuals}
\label{do:sec:cf}

Pearl's third rung asks what $Y$ would have been had $X$ been $x$, in the world that actually
occurred. In a structural model this is abduction (infer $U$ from the evidence), action
(mutilate), prediction (compute).

\begin{definition}[Exogenous record]
\label{do:def:record}
A \emph{run} of $\sem{M}$ is a causal history together with, at each gadget, the candidate that
fired. Its \emph{exogenous record} $u$ is the vector of fired ancilla classes. A run is
\emph{replayed under} $\dok(C)$ from a pre-state $S$ by executing $\dok(C)\ \ink\ S$ with the
resolver instructed, at each gadget whose contention component is unchanged from the recorded run,
to fire the recorded class if it is still a candidate of nonzero restricted weight, and otherwise
to sample.
\end{definition}

\begin{theorem}[Counterfactuals are replay]
\label{do:thm:cf}
Let $M$ be positive, $u$ the exogenous record of a run with outcome $(v_1,\dots,v_n)$, and $S$ the
initial state. Then replaying under $\dok(u^{(i)}_1 = x)$ from $S$ yields, at every gadget other
than $i$ and its descendants, the recorded value, and at $i$ and its descendants the values
$Y_{V_i = x}(u)$ of the structural model whose exogenous variables are the ancilla draws. On the
staged fragment the ``otherwise sample'' branch is never taken.
\end{theorem}

\begin{proof}
By Lemma~\ref{do:lem:staging} every gadget's contention component is itself, and its candidate set
depends only on which parent values arrived, every ancilla value being present. So the recorded
class is always a candidate; it has nonzero restricted weight by positivity except at gadget $i$,
where the restriction forces $x$. Non-descendants of $V_i$ receive the same parent values as in the
run and fire the recorded class; descendants receive the mutilated values and fire their recorded
class, which is $f_j(\mathit{pa}_j', u_j)$ for the structural model in which $u_j$ selects the
class --- the standard reduction of a CBN to a structural model by inverse-CDF sampling, with the
class index as $U_j$.
\end{proof}

The ``otherwise sample'' clause is where the general case lives: off the staged fragment an
intervention can change what rests on a channel and thereby which candidates exist, and ``the same
draw'' is no longer well defined. That the record of a run suffices for its counterfactuals is also
the sense in which the history monad of Chapter~\ref{ch:histmonad} was always the third rung
waiting to be read as one: $\Hist$ records; replay under a $\dok$ is the use of the record.

\subsection{The dictionary}

\begin{table}[t]
\centering\small
\begin{tabular}{@{}>{\raggedright\arraybackslash}p{0.30\textwidth}>{\raggedright\arraybackslash}p{0.62\textwidth}@{}}
\toprule
\textbf{Pearl} & \textbf{graded rho over $\Rnn$}\\
\midrule
variable $V_i$ & gadget $i$; value $=$ fired ancilla class\\
parents $\mathrm{PA}_i$ & data on the channels the pattern and clause read\\
mechanism $P_i(\cdot\mid\mathit{pa})$ & the clause, in disjunctive normal form\\
exogenous $U_i$ & the resolver's draw at the ancilla race\\
DAG $G$ & causal order of the staged term\\
$\Pr(y \mid x)$ & post-selection on histories\\
$\Pr(y \mid \dok(x))$ & $\dok(u_1 = x)\ \ink\ \sem{M}$\\
$\sigma$-intervention & $\dok(y \models \psi)$\\
$Y_x(u)$ & replay under $\dok$ with the record held\\
positivity & no restricted receipt starves\\
d-separation & a crisp-sort formula over the staged term (conjectured)\\
\bottomrule
\end{tabular}
\caption{The dictionary.}
\label{do:tab:dictionary}
\end{table}

Table~\ref{do:tab:dictionary} collects the correspondence. Two entries deserve a word.
\emph{Set-valued and graded constraints} --- $\dok(y \in S)$ with $|S| \geq 2$, or
$\dok(y \models \psi)$ with $\psi$ graded --- are not Pearl's atomic $\dok$ but the
$\sigma$-interventions of Correa and Bareinboim~\cite{CorreaBareinboim}, which replace a
mechanism by an arbitrary new one and state their rules on the graph whose incoming edges at $X$
are those of the \emph{new} mechanism. Here the construct is natively a $\sigma$-intervention,
and the new mechanism's parent set is read off the free names of the constraint: a $\dok$ whose
constraint reads other variables is a policy, and its parents are written in the formula. The
induced graph is syntactic, for the intervention as for the model. \emph{Conditioning} is
post-selection on histories, which costs $1/\Pr(x)$ retries; the chapter on the complex instance
met the same quantity as the price of a herald, and the ledger charges it the same way.

\section{What it costs, and the fragment that tames it}
\label{do:sec:complexity}

The membrane semantics is simple to state and expensive to run, and the expense has three
independent sources. The constraint language can contain modalities and fixed points, which cost
what Chapter~\ref{ch:grading} said they cost: a clause that looks at successors is not local. The
scope can be dynamic through reflection, in which case the $\dok$ is a run-time membrane and not a
rewrite. And restriction changes the normalization of contention components whose enumeration is
permanent-shaped (Remark~\ref{gw:rmk:costs}). A fourth question --- whether a given $\dok$ is
vacuous, fatal, or neither --- is reachability, hence undecidable in general.

The response is to grade the constraint language into fragments by evaluation cost and to impose
type disciplines on the scope, and the disciplines are not new to this book. \emph{Staged}: every
receipt is its own contention component --- the quantum-safe fragment of Chapter~\ref{sec:quantum}.
\emph{Flat}: no received code can come to contain a receipt --- the hypothesis of
Proposition~\ref{do:prop:lowering}. \emph{Finite-domain}: channel types with finitely many values
--- what a language definition supplies for free. \emph{Confined}: a $\dok$ reaches only code its
author has paid to deploy --- the seal discipline of Chapter~\ref{ch:carho}.

\begin{proposition}[The practical fragment]
\label{do:prop:practical}
On staged, flat, finite-domain, confined terms with constraints that are indicators of finite sets:
(i) $\dok(C)\ \ink\ P$ lowers at elaboration and the membrane is erased; (ii) each constrained
candidate costs one structural-congruence test per element of the set; (iii) vacuity and fatality
are decidable; (iv) the induced graph is the static DAG with the constrained variables' incoming
edges deleted; and (v) the semantics is that of Definition~\ref{do:def:cbn} exactly. With
constraints of bounded modal depth instead, (i)--(iii) hold and the semantics is a
$\sigma$-intervention on a CBN.
\end{proposition}

\begin{proof}
(i) is Proposition~\ref{do:prop:lowering}, whose hypothesis is flatness. (ii) is the definition
of the constraint fragment. (iii): under staging every constrained receipt fires exactly once and
its candidate set is determined by its parents' values, which range over a finite product;
enumerate. (iv) and (v) are Theorem~\ref{do:thm:surgery}, whose hypotheses staging and finite
domains supply.
\end{proof}

That four disciplines introduced for four purposes --- a quantum state, an exact lowering, a
decidable check, an owner's seal --- draw the same line is the best evidence we have that the line
is real. Outside it nothing fails; Pearl's calculus stops and the calculus of this book continues.
Off staging, the contention components are the branching cells of a confused event
structure~\cite{AbbesBenveniste} and a $\dok$ that touches one is a restriction of the cell's
distribution; off flatness, a $\dok$ is a run-time membrane; off finite domains, positivity is a
reachability question; off confinement, an intervention on another's code is a soft intervention
by Proposition~\ref{do:prop:nocontext} --- an experiment rather than a surgery, which is what
\chSci's scientist does to its environment all day.

\section{Execution, briefly}
\label{do:sec:execution}

The construct was designed against the execution framework of f1r3node rather than for a
simulator, and the book returns to that framework in its own right when it turns from the calculus
to the machine. Four facts are worth recording here so that the reader knows the construct is not
only a semantics. Lowering is one more rewrite in the elaboration pass that already normalizes
every deployment, and leaves no run-time trace. Off the flat fragment the $\dok$ is a run-time
membrane and the twin of the cost membrane: a scoped decoration on the cut, read off the term by
the reducer, altering the COMM rule within; the two commute, since metering does not read the
clause and the clause does not read the meter. An intervention is priced as one COMM plus the
marginal cost of its conjunct, by Principle~\ref{do:prin:metered} --- though fixing the charge
against the ledger is owed (\S\ref{do:sec:open}). And a counterfactual needs a pre-state and a
record of fired candidates, which is exactly what a validator that executes a block and records
the winners already has; the losing candidates are recomputable from the pre-state, so the same
instrumentation that serves counterfactual replay serves the gradient of
Theorem~\ref{sci:thm:score}, one number per firing.

One observation from the execution model belongs here rather than later, because it is the
dictionary's last row. A shard in which one node executes and the others replay to confirm the
state change is an experimenter/spectator split: the executing node resolves the races and so
stands on rung two; the replayers observe and stand on rung one. $\Pr(y \mid x)$ estimated from
the finalized trace is the trace distribution and no more. That a distributed ledger is, by its
design, a machine that separates the two rungs is not something its designers set out to build.

\section{What this does for the scientist}
\label{do:sec:scientist}

\chSci\ will build a learner whose hypotheses are graded clauses and whose experiments are guards.
Everything it does is, in Pearl's vocabulary, rung one and a half: it observes, and it perturbs
its environment by composing contexts with it, and by Proposition~\ref{do:prop:nocontext} a
context perturbs without setting. The $\dok$ is what gives it rung two over the code it owns. A
hypothesis is a clause; a clause under a $\dok$ is a hypothesis together with the experiment that
would test it; and --- this is the point the reproduction section of that chapter will make
exact --- a germ carrying both is a hedge that is inherited, since a $\dok$ is a term and
crossover on quoted code applies to it unchanged. The construct of this chapter is what the
scientist's germs were missing in order to be rung-two objects rather than rung-one ones.

The Prestige takes this up as the constructive half of an argument it makes three times: that a
system whose witnesses reduce confluently can \emph{state} a non-confluent phenomenon but cannot
\emph{inhabit} it, so that a proof assistant of the usual kind, and a transformer trained on
completed runs, both stand on rung one by construction. Rung two is a graded race under a $\dok$;
rung three is its replay; and the smallest behavior invisible to a trace --- Milner's two vending
machines, which agree on every completed run and differ in a \emph{refusal} --- becomes a one-line
program that, under $\dok(y = \mathit{tea})$, either delivers tea or starves. That chapter is
where the argument is made. This one is where the thing it leans on is built.

\section{What is owed}
\label{do:sec:open}

\begin{enumerate}[leftmargin=1.6em,itemsep=3pt]
\item \textbf{d-separation as a behavioral formula.} Corollary~\ref{do:cor:rules} evaluates its
premises on $G$. Since $G$ is the causal order of the staged term, the premises should be
crisp-sort $\nu$-properties over $\sem{M}$ --- ``no history has a causal path from a $Z$-firing to
a $Y$-firing not passing through $X$ or $W$'' --- dischargeable by the model checker. Stated, not
proved.
\item \textbf{The rules off the staged fragment.} Outside staging the contention components are
branching cells and the event structure may be confused. The conjecture is that the
$\sigma$-calculus rules hold on the confusion-free fragment with d-separation read off the event
structure, and on the confused fragment per cell with the cell's distribution in place of the
mechanism.
\item \textbf{The price of a $\dok$.} Principle~\ref{do:prin:metered} says an intervention is a
COMM. The marginal charge of the constraint conjunct, the budget charge of a modal constraint, and
the charge for a counterfactual replay have not been fixed against the ledger of
Chapter~\ref{ch:carho}. Until they are, ``an intervention costs one token'' names an intention.
\item \textbf{Identifiability across namespaces.} Pearl's question --- when is $\Pr(y \mid \dok(x))$
computable from observational data in the presence of hidden variables --- has a sharp rendering:
the hidden variables are the code one cannot read, the boundary is the namespace
(Chapter~\ref{ch:scope}), and the observational data is the finalized trace. The identification
algorithm~\cite{ShpitserPearl} should transfer to the staged fragment directly; what it says
about a shard is the question.
\item \textbf{Continuous domains and cycles.} The gadget enumerates the domain on the ancilla
channel; a continuous variable needs an ancilla carried as a density and the normalization become
an integral, for which the Gillespie reading of Chapter~\ref{ch:grading} is the natural home. A
recursive term has an acyclic event-level graph and a cyclic variable-level one; the cyclic
structural models of the literature and the graded fixed point should meet there, with the budget
as the unrolling depth.
\end{enumerate}

% ===================================================================
